% =========================================================
% Volume IV  \S{}1.7.?  --- R as a Model of the Ordered Field Signature
% WIRING NOTE (cross-volume):
%   Verify exact labels before final commit.
%
% Expected targets:
%   def:arithmetic-signature
%   def:algebraic-field
%   def:positive-cone
%   def:least-upper-bound-property
%   thm:r-is-field
%   thm:r-is-ordered-field
%   ax:completeness-of-the-real-numbers
%   thm:q-embeds-in-r
%   thm:r-contains-q
%   def:structure-expansion
%   def:structure-reduction
% =========================================================

\subsection*{The Real Numbers as a Structure}

\begin{remark*}[Modeling Goal]
We exhibit $\mathbb{R}$ as a concrete model of the ordered-field signature.
The field structure is first-order and is captured by satisfaction of the
ordered-field axioms. The completeness of $\mathbb{R}$ is then discussed
separately, because it does not belong to the ordinary first-order field
signature.
\end{remark*}

% ---------------------------------------------------------
% Signature
% ---------------------------------------------------------

\begin{remark*}[Exposition]

A model consists of a signature together with an interpretation on a carrier.

The signature considered here is the ordered-field signature

\[
\mathcal L_{\mathrm{of}}
=
\{+,\cdot,-,{}^{-1},0,1,<\}.
\]

It contains the field operations together with a distinguished order relation.

\end{remark*}

% ---------------------------------------------------------
% Model
% ---------------------------------------------------------

\begin{definition}[The Real Number Model $\mathcal R$]
\label{def:real-number-model}

The \textbf{real number model} is the
$\mathcal L_{\mathrm{of}}$-structure

\[
\mathcal R
=
\bigl(
\mathbb R;
+^{\mathcal R},
\cdot^{\mathcal R},
-^{\mathcal R},
{}^{-1\,\mathcal R},
0^{\mathcal R},
1^{\mathcal R},
<^{\mathcal R}
\bigr),
\]

whose carrier

\[
|\mathcal R|
=
\mathbb R
\]

is the set of real numbers constructed in Volume~II, and whose symbols are
interpreted by the real-number operations and order.

\end{definition}

% ---------------------------------------------------------
% Quantified statement
% ---------------------------------------------------------

\begin{remark*}[Standard quantified statement]

\[
|\mathcal R|
=
\mathbb R,
\]

\[
+^{\mathcal R},
\cdot^{\mathcal R}
:
\mathbb R\times\mathbb R
\to
\mathbb R,
\]

\[
-^{\mathcal R}
:
\mathbb R
\to
\mathbb R,
\]

\[
{}^{-1\,\mathcal R}
:
\mathbb R\setminus\{0\}
\to
\mathbb R,
\]

\[
0^{\mathcal R},
1^{\mathcal R}
\in
\mathbb R,
\]

\[
<^{\mathcal R}
\subseteq
\mathbb R\times\mathbb R.
\]

\end{remark*}

% ---------------------------------------------------------
% Dependencies
% ---------------------------------------------------------% ---------------------------------------------------------
% Interpretation
% ---------------------------------------------------------

\begin{remark*}[Interpretation]

The structure

\[
\mathcal R
\]

packages the entire real-number system into a single mathematical object.

The carrier, operations, inverses, and order already exist from the Volume~II
construction. The model simply presents them as a unified structure that may
serve as an argument to general theorems about ordered fields.

\end{remark*}

% ---------------------------------------------------------
% Satisfaction
% ---------------------------------------------------------

\begin{remark*}[Satisfaction]

The structure

\[
\mathcal R
\]

is a model of the ordered-field axioms:

\[
\mathcal R
\models
\Sigma_{\mathrm{of}},
\]

equivalently,

\[
\operatorname{OrderedField}
(\mathbb R,+,\cdot,<,0,1).
\]

This is witnessed by

\hyperref[thm:r-is-ordered-field]{$\mathbb R$ Is an Ordered Field}

in Volume~II.

In particular,

\[
\mathcal R
\models
\Sigma_{\mathrm{field}}
\]

because

\hyperref[thm:r-is-field]{$\mathbb R$ Is a Field}

verifies the field axioms.

\end{remark*}

% ---------------------------------------------------------
% Completeness discussion
% ---------------------------------------------------------

\begin{remark*}[Completeness Is Separate]

The ordered-field structure does \emph{not} yet distinguish
$\mathbb R$ from every other ordered field.

The defining analytic property of the real numbers is completeness, expressed
in Volume~II by

\[
LeastUpperBoundProperty(\mathbb R).
\]

Equivalently, every nonempty subset of $\mathbb R$ that is bounded above
possesses a least upper bound.

This property is witnessed by

\hyperref[ax:completeness-of-the-real-numbers]
{Completeness of the Real Numbers}.

\end{remark*}

\begin{remark*}[Important Structural Distinction]

The statement

\[
\mathcal R
\models
\Sigma_{\mathrm{of}}
\]

belongs to the first-order ordered-field structure.

The completeness property does not belong to the ordinary ordered-field
signature.

Consequently the real-number model should be viewed as

\[
\text{ordered field}
+
\text{completeness}.
\]

These are conceptually distinct layers.

\end{remark*}

% ---------------------------------------------------------
% Reduction to field
% ---------------------------------------------------------

\begin{remark*}[Reduction to the Field Structure]

Forgetting the order relation leaves the reduct

\[
\mathcal R
\!\restriction\!
\{+,\cdot,-,{}^{-1},0,1\},
\]

which is simply the field structure of the real numbers.

\[
\operatorname{Reduction}
\Bigl(
\mathcal R
\!\restriction\!
\{+,\cdot,-,{}^{-1},0,1\},
\mathcal R
\Bigr).
\]

\end{remark*}

% ---------------------------------------------------------
% Additive reduct
% ---------------------------------------------------------

\begin{remark*}[Reduction to the Additive Group]

Forgetting multiplication and order leaves

\[
(\mathbb R;+,-,0),
\]

the additive abelian group of real numbers.

\end{remark*}

% ---------------------------------------------------------
% Blueprint position
% ---------------------------------------------------------

\begin{remark*}[Blueprint Position]

The real-number model occupies the final stage of the standard ordered
number-system hierarchy:

\[
\mathbb N
\hookrightarrow
\mathbb Z
\hookrightarrow
\mathbb Q
\hookrightarrow
\mathbb R.
\]

The final embedding is witnessed by

\hyperref[thm:q-embeds-in-r]
{$\mathbb Q$ Embeds in $\mathbb R$}.

The rationals appear inside $\mathbb R$ as a subfield.

\end{remark*}

% ---------------------------------------------------------
% Interpretation
% ---------------------------------------------------------

\begin{remark*}[Interpretation]

The passage

\[
\mathbb Q
\longrightarrow
\mathbb R
\]

does not add new field operations.

Addition, subtraction, multiplication, and division already exist in
$\mathbb Q$.

Instead it adds completeness.

This closes the gaps left by the rational numbers and supplies limits,
suprema, infima, and the analytic foundation required for calculus and
real analysis.

\end{remark*}

% ---------------------------------------------------------
% Structural significance
% ---------------------------------------------------------

\begin{remark*}[Structural Significance]

The real-number model is simultaneously

\begin{itemize}
    \item a field,
    \item an ordered field,
    \item a complete ordered field.
\end{itemize}

The first two properties belong to its algebraic structure.

The third is the property that makes real analysis possible.

For this reason $\mathbb R$ occupies the analytic layer of the structural
hierarchy and serves as the foundational model for limits, continuity,
differentiation, and integration.

\end{remark*}

\begin{dependencies}

\begin{itemize}
    \item \hyperref[def:algebraic-field]{Field}.
    \item \hyperref[def:positive-cone]{Positive Cone}.
    \item \hyperref[def:reals]{The Real Numbers}.
    \item \hyperref[thm:r-is-a-field]{$\mathbb R$ Is a Field}.
    \item \hyperref[thm:r-is-an-ordered-field]{$\mathbb R$ Is an Ordered Field}.
    \item \hyperref[thm:r-is-complete-ordered-field]{$\mathbb R$ Is a Complete Ordered Field}.
\end{itemize}

\end{dependencies}
